\documentclass[UTF8,a4paper,12pt,fontset=fandol]{ctexart}

\usepackage{geometry}
\usepackage{graphicx}
\usepackage{booktabs}
\usepackage{array}
\usepackage{enumitem}
\usepackage{hyperref}
\usepackage{xcolor}

\geometry{left=2.5cm,right=2.5cm,top=2.7cm,bottom=2.7cm}

\hypersetup{
  colorlinks=true,
  linkcolor=blue!60!black,
  citecolor=blue!60!black,
  urlcolor=blue!60!black
}

\setlist[itemize]{itemsep=0.25em,topsep=0.25em}
\setlist[enumerate]{itemsep=0.25em,topsep=0.25em}

\title{\textbf{多轮对话 Agent 推理系统优化小组报告}}
\author{
  组员 A \quad 组员 B \quad 组员 C \quad 组员 D\\
  \small 指导教师：XXX
}
\date{\today}

\begin{document}

\maketitle

\begin{abstract}
本文围绕多轮对话 Agent 推理系统的首 token 延迟（Time To First Token, TTFT）优化展开分析。报告首先说明小组分工与问题拆解方式，然后从请求排队、prefill、KV cache、上下文长度和调度策略等维度分析 TTFT 的主要影响因素，进一步总结若干尝试过的优化策略，最后给出综合权衡后的最终方案。
\end{abstract}

\tableofcontents
\newpage

\section{项目背景与目标}

多轮对话 Agent 推理系统通常需要在连续交互中理解用户意图、维护上下文状态、调用外部工具并生成最终答案。相关研究包括链式思维提示、推理与行动结合的 ReAct 框架，以及面向工具使用的语言模型方法\cite{wei2022chain,yao2023react,schick2023toolformer}。

本项目关注的核心性能指标是首 token 延迟（Time To First Token, TTFT），即从用户请求到模型开始输出第一个 token 的时间。对于交互式 Agent 系统，TTFT 直接影响用户对系统是否“响应及时”的主观感受，也能反映请求进入模型前后的关键链路开销。

本项目的目标是围绕 TTFT 对多轮对话 Agent 推理流程进行系统性分析与优化。重点关注以下问题：

\begin{itemize}
  \item 如何拆解一次 Agent 请求的 TTFT 组成；
  \item 如何分析上下文长度、KV cache 命中情况和 prefill 计算对 TTFT 的影响；
  \item 如何减少多轮对话中重复上下文带来的首 token 等待时间；
  \item 如何建立可复现实验流程，用 TTFT 数据评估不同优化策略的效果。
\end{itemize}

\section{组员分工}

\begin{table}[htbp]
  \centering
  \caption{小组成员分工}
  \begin{tabular}{p{3cm}p{5cm}p{6cm}}
    \toprule
    成员 & 主要职责 & 具体工作 \\
    \mid
    组员 A & 系统分析与方案设计 & 性能瓶颈分析，设计基于cache aware的调度方案以及衍生cache策略 \\
    组员 B & 算子融合 &  \\
    组员 C & 超参数分析 &  \\
    \bottomrule
  \end{tabular}
\end{table}

\section{问题拆解}

我们将一次多轮对话 Agent 请求的 TTFT 拆解为请求进入系统、上下文构造、模型排队、prefill 计算、KV cache 读取或写入，以及首 token decode 六个阶段。不同阶段的耗时共同决定用户看到第一个 token 前的等待时间。

\subsection{请求进入系统}

请求进入系统后通常会经历鉴权、路由、会话读取和任务类型判断。若 Agent 还需要先执行工具调用或规划步骤，则这些前置操作也会被计入首 token 之前的等待时间。

\subsection{上下文构造}

上下文构造阶段需要拼接系统提示词、历史对话、工具结果和当前用户输入。多轮场景下，历史消息越长，prefill 需要处理的 token 越多，TTFT 通常也越高。

\subsection{模型排队}

当推理服务负载较高时，请求可能先进入队列等待 batch 调度。队列等待不属于模型计算本身，但会直接增加 TTFT，并且容易在高并发场景中成为主要波动来源。

\subsection{Prefill 计算}

prefill 阶段会对输入 token 做一次前向计算，为后续 decode 生成注意力状态。对于首 token 来说，prefill 往往是最关键的模型侧开销，尤其在长上下文或 batch 较大时更加明显。

\subsection{KV Cache}

KV cache 保存历史 token 的 key/value 状态，使后续 token 不必重复计算完整上下文。在多轮对话中，如果系统能够复用前几轮稳定不变的上下文缓存，就可以减少 prefill 计算量，从而降低 TTFT。反之，如果每轮都重新拼接并计算完整历史，KV cache 的收益就无法充分体现。

\subsection{首 Token Decode}

首 token decode 是 prefill 完成后的第一次生成步骤。与长文本生成的总耗时相比，首 token decode 通常不是主要瓶颈，但它仍会受到采样参数、模型规模和服务端实现的影响。

\section{性能分析}

\subsection{核心指标}

本项目只关注一个核心性能指标：\textbf{TTFT}。其定义为用户请求到达 Agent 服务后，到客户端接收到模型输出第一个 token 为止的时间。为避免统计口径混乱，我们不将总生成时长、吞吐量、token 成本或任务成功率作为主要性能指标。

实验中记录的 TTFT 可以进一步拆分为：

\begin{itemize}
  \item \textbf{端到端 TTFT}：从请求进入 Agent 服务到首 token 返回；
  \item \textbf{模型侧 TTFT}：从请求进入推理后端到首 token 生成；
  \item \textbf{排队时间}：请求等待调度和 batch 合并的时间；
  \item \textbf{prefill 时间}：模型处理输入上下文并建立 KV 状态的时间；
  \item \textbf{首 token decode 时间}：完成 prefill 后生成第一个 token 的时间。
\end{itemize}

\subsection{影响因素}

初步分析表明，TTFT 主要受到以下因素影响：

\begin{enumerate}
  \item \textbf{输入上下文长度}：上下文越长，prefill 计算量越大；
  \item \textbf{KV cache 命中率}：可复用的历史上下文越多，重复 prefill 越少；
  \item \textbf{请求调度策略}：batch 大小、排队策略和优先级会影响首 token 等待时间；
  \item \textbf{工具前置调用}：若 Agent 在模型输出前必须调用工具，会拉长端到端 TTFT；
  \item \textbf{模型与硬件配置}：模型规模、量化方式、GPU 显存压力和并发负载都会影响模型侧 TTFT。
\end{enumerate}

\subsection{KV Cache 分析}

在多轮对话中，前几轮对话、系统提示词和工具说明通常在相邻请求之间保持不变。这部分内容适合通过 KV cache 复用，避免每轮重新 prefill。分析 KV cache 时需要重点关注以下问题：

\begin{itemize}
  \item 哪些上下文片段在多轮请求中保持稳定；
  \item 当前推理框架是否支持 prefix cache 或 session cache；
  \item prompt 拼接方式是否会破坏缓存前缀一致性；
  \item cache 命中后 prefill 时间下降是否能体现在端到端 TTFT 中；
  \item 高并发时 cache 占用显存是否会反过来影响调度和排队。
\end{itemize}

\subsection{主要瓶颈}

结合 TTFT 统计和链路拆解，我们认为主要瓶颈集中在三个方面：第一，历史上下文不断增长导致 prefill 时间上升；第二，prompt 拼接不稳定导致 KV cache 难以命中；第三，高并发下排队时间波动较大，使用户感知到的首 token 延迟不稳定。

\section{尝试策略}

\subsection{策略一：上下文裁剪}

对历史对话进行裁剪，只保留与当前轮请求直接相关的内容。该策略可以降低 prefill token 数，从而减少 TTFT，但需要避免删掉影响任务理解的关键信息。

\subsection{策略二：稳定 Prompt 前缀}

将系统提示词、工具说明和固定格式模板放在稳定前缀中，减少每轮请求中前缀内容的变化。稳定前缀有利于 prefix KV cache 命中，适合多轮对话 Agent 的连续请求场景。

\subsection{策略三：会话级 KV Cache 复用}

为同一个会话维护可复用的 KV cache，使新一轮请求只对新增用户输入和必要增量上下文进行 prefill。该策略对降低多轮 TTFT 最直接，但需要管理 cache 生命周期、显存占用和会话淘汰策略。

\subsection{策略四：工具调用后置或异步化}

对于不影响首轮响应的工具调用，尝试后置或异步执行，先让模型返回初始响应或状态提示。该策略可以降低端到端 TTFT，但只适用于工具结果不是首 token 必需条件的场景。

\subsection{策略五：请求调度优化}

通过限制最大 batch 等待时间、区分短上下文和长上下文请求、为交互式请求设置更高优先级等方式降低排队时间。该策略主要缓解高并发下 TTFT 抖动。

\section{最终策略}

综合实验结果与实现成本，我们采用“稳定 prompt 前缀 + 上下文裁剪 + 会话级 KV cache 复用 + 交互式请求优先调度”的组合策略。该方案直接围绕 TTFT 的主要组成部分进行优化：减少 prefill token 数、提高 cache 命中率，并降低排队时间。

\subsection{策略流程}

\begin{enumerate}
  \item \textbf{固定前缀构造}：将系统提示词、工具 schema 和固定任务模板保持稳定，避免无意义的顺序或格式变化；
  \item \textbf{上下文增量更新}：每轮只追加新的用户输入、必要工具结果和压缩后的历史摘要；
  \item \textbf{KV cache 复用}：对稳定前缀和已确认历史建立会话级缓存，新请求优先复用已有 KV 状态；
  \item \textbf{请求分级调度}：将交互式 Agent 请求与后台批处理请求分开排队，减少首 token 等待抖动；
  \item \textbf{TTFT 监控}：分别记录端到端 TTFT、排队时间、prefill 时间和 cache 命中情况，便于定位回归。
\end{enumerate}

\subsection{预期收益}

该组合方案能够在保持 Agent 多轮上下文能力的同时降低首 token 等待时间。上下文裁剪减少了必须计算的 token 数，稳定 prompt 前缀和会话级 KV cache 提升了缓存复用率，请求分级调度则减少了高并发下的排队波动。最终优化目标不是减少总生成时间，而是让用户更快看到第一个 token。

\section{总结与展望}

本报告围绕多轮对话 Agent 推理系统的 TTFT 优化展开，将首 token 延迟拆解为请求处理、上下文构造、模型排队、prefill、KV cache 和首 token decode 等阶段。相比泛泛讨论准确率、成本或总响应时长，聚焦 TTFT 能更直接地指导交互式 Agent 的性能优化。

最终方案强调三点：保持 prompt 前缀稳定，尽可能复用 KV cache；控制多轮上下文长度，减少重复 prefill；优化请求调度，降低高并发下的排队波动。后续工作可以进一步补充真实压测数据和消融实验，比较不同上下文长度、cache 命中率和并发负载下的 TTFT 变化。


\bibliographystyle{plain}
\bibliography{references}

\end{document}
